\section{Dimension-scaled costs of the finite-phase separation}\label{sec:dimension-cost}
The exact phase identity permits an explicit succession of finite pairs, rather than unspecified successive perturbation limits at each fixed dimension. This is a quantitative refinement of the preceding family, not a new preparation scheme or an optimal resource frontier. The noise parameter in this section changes with dimension.

\begin{theorem}[An achievable information--conditioning tradeoff]\label{thm:dimension-cost}
Fix $t>0$ and $\beta>1$ independently of the integer $d\ge2$, and put $\nu=\beta+3>4$. In Theorem~\ref{thm:phase} choose
\[
 e^2=t/d,\quad \theta=1-d^{-\beta},\quad \phi=\pi,
 \quad\lambda_d=-1/(d^2-1).
\]
The faithful input pair has reference dimension $d$, the stated norm-one local Hamiltonian, and exact actual-marginal order
\[
 K_d=d^2-d^{2-\beta}+d^{-\beta}.
\]
Write $I_d=D_{\mathrm U}(\rho_\pi\Vert s_\theta)$, let $R_d$ be its output/input relative-entropy retention, and let $\operatorname{cond}(s_\theta)$ be the ratio of largest to smallest eigenvalue. Then
\begin{align}
 dR_d&\longrightarrow\alpha=\frac{\nu-4}{(1+t)\nu},\label{eq:scaled-retention}\\
 \frac{d^2I_d}{\log d}&\longrightarrow4t\nu,\qquad
 \frac{\operatorname{cond}(s_\theta)}{d^\nu}\longrightarrow\frac1t,\label{eq:scaled-cost}\\
 \frac d2\|\rho_\pi-s_\theta\|_1&\longrightarrow2\sqrt t.
\end{align}
In particular $R_d/b_{d^2}(\lambda_d)\sim\alpha d$, although both $R_d$ and the absolute input/output information tend to zero. The information of the output is asymptotic to $4t(\nu-4)\log(d)/[(1+t)d^3]$.
\end{theorem}
\begin{proof}
Set $x=e^2$, $\delta=1-\theta$, $q=d-1+x$ and $\Xi=d(d-1)-1+(d-1)x$. Directly from the subnormalized $W_p$,
\begin{align*}
 T_{\rm in}&=1-\frac{\delta\Xi}{dq},&
 \Delta_{\rm in}&=\frac{x\delta[d^2-(d^2-1)\delta]}{d^2q^2},\\
 T_{\rm out}&=\frac{d[1+(d-1)x]+\delta\Xi}{(d^2-1)dq},&
 \Delta_{\rm out}&=\frac{x\delta(d^2-\delta)}{(d^2-1)d^2q^2},
\end{align*}
where $T$ denotes trace and $\Delta$ determinant, not divergence. Substituting $x=t/d$, $\delta=d^{-\beta}$ gives
\[
 T_{\rm in}\to1,\quad d^\nu\Delta_{\rm in}\to t,
 \quad d^3T_{\rm out}\to1+t,\quad d^{\nu+2}\Delta_{\rm out}\to t.
\]
These limits also hold for noninteger fixed $\beta>1$: use $u=1/d$ and $y=d\delta$, and evaluate the corresponding rational functions jointly at $(u,y)=(0,0)$. Their denominators are nonzero there for fixed $t>0$.

For a positive two-dimensional matrix with trace $T$ and determinant $\Delta$, if $\Delta/T^2\to0$ then $w_+\sim T$ and $w_-\sim\Delta/T$. Therefore
\[
 \begin{array}{c|cc}
 &w_+&w_-\\\hline
 \mathrm{input}&1+o(1)&t d^{-\nu}(1+o(1))\\
 \mathrm{output}&(1+t)d^{-3}(1+o(1))&
       \dfrac{t}{1+t}d^{-(\nu-1)}(1+o(1)).
 \end{array}
\]
It follows that $\ell(W_\theta)/\log d\to\nu$ and
$\ell(W_{\lambda_d\theta})/(d^3\log d)\to(\nu-4)/(1+t)$. Since $d^4\lambda_d^2\to1$, the exact finite-phase identity proves~\eqref{eq:scaled-retention}. Its physical coherence obeys $d^2c^2\to t$, giving $I_d=4c^2\ell(W_\theta)$ and the information limit. The trace distance is exactly $2|c|$.

The exact minimum eigenvalue of the full input state is
\[
 \lambda_{\min}(s_\theta)=\frac{\delta\min(1,x)}{dq}.
\]
Indeed $s_\theta\succeq\delta B$, and an off-Schmidt coordinate with a reference label of minimum weight attains equality. Eventually $x=t/d\le1$, so $d^\nu\lambda_{\min}\to t$. The largest eigenvalue tends to one, since it lies between $\theta$ and one. This proves the condition-number limit, including all ambient modes. The order and faithful positivity follow from Theorem~\ref{thm:phase} at each finite dimension, before taking limits.
\end{proof}

For any fixed $0<\gamma<1$, choose $t>0$ with $\gamma(1+t)<1$ and then
\[
 \nu>\frac4{1-\gamma(1+t)},\qquad\beta=\nu-3.
\]
This yields $\alpha>\gamma$ with parameters independent of $d$. The conditioning exponent and constants depend on $\gamma$; no fixed exponent works uniformly as $\gamma\uparrow1$, and $\alpha=1$ is not attained with finite parameters. This approximates the scale of an existing lower, not an unrestricted optimal coefficient.

\paragraph{A different order cost and necessary quantifiers.}
The alternative explicit choice $e=1/d$, $\theta=1-1/d$ gives $dR_d\to1/10$, $I_d\sim20\log(d)/d^3$, condition number $\sim d^5$ and trace distance $\sim2d^{-3/2}$. Its exact order is $d^2-d+1/d$. The choice $t=1,\beta=2$ above retains the same leading retention and condition exponent with input information larger by a factor asymptotic to $d$, but spends the larger order $d^2-1+d^{-2}$. It is not a fixed-budget improvement.

Fixed parameters are substantive. If $t$ is instead replaced by $t_d=1/d$ with $\beta=2$, the true limit is $dR_d\to1/3$, not the $1/5$ obtained by inserting $t_d$ into the fixed-$t$ formula and then taking $d\to\infty$. The input information is then $24\log(d)/d^3$, not the analogous constant20. Likewise $\beta=1$ is outside this theorem: it changes the output-root regime. No uniform substitution of variable parameters is justified.

\paragraph{Evidence and operational boundary.}
Audit1068 reconstructed the rational and logarithmic limits, including noninteger $\beta$, and checked seven full ambient constructions,57 integer-interval cases and containment of the40 supplied interval certificates. For the $e=1/d$, $\theta=1-1/d$ sequence, it retains negative gaps versus Bell at $d=2,\ldots,7$; positive checked points beginning at8 do not prove a threshold for every integer $d\ge8$. Nine of the32 original fixed tradeoff points are also adverse. The limit proves eventual separation, without assigning a universal finite threshold from the table. Initial control failures and the correction of a matrix-indexing error remain documented.

This is a relative separation under $\lambda_d\to0$, approaching the replacement channel, not a growing absolute advantage at fixed noise. If $d=2^m$, the condition number is exponential in $m$ and the input information is exponentially small up to a factor $m$. Locality means $H_A\otimes I_R$, not a proved few-body implementation on $m$ qubits. The SLD information of the same input phase is $4c^2/T_{\rm in}\sim4t/d^2$ and does not grow. No preparation complexity, experimental tolerance, Pareto frontier, necessary resource exponent, optimal absolute estimation or MI/capacity consequence follows. The specific quantitative scaling is a direct refinement once the preceding exact block identity is available; its priority is not certified.
